\chapter{Data Preparation}
\label{ch:data_preparation}

Questa fase ha lo scopo di organizzare, pulire e trasformare i dati grezzi estratti precedentemente. L'obiettivo è ottenere un dataset formattato in modo ottimale per le fasi successive di addestramento dei modelli di Intelligenza Artificiale.

% ─────────────────────────────────────────────────────────────────────────────
\section{Modello Red Flag Detection (NLP)}
\label{sec:prep_red_flag}

A seguito dell'esplorazione del dataset (come illustrato nel Capitolo~\ref{ch:data_understanding}), abbiamo messo a punto uno script Python per ripulire i testi. Le decisioni prese in questa fase sono fondamentali: dobbiamo assicurarci di rimuovere tutto il "rumore" del web senza però cancellare le sfumature cliniche ed emotive che potrebbero rivelare un rischio suicidario.

Di seguito analizziamo i vari passaggi eseguiti, disponibili e riproducibili all'interno del Notebook di analisi \texttt{02\_RedFlag\_DataPreparation.ipynb}.

\subsection{1. Rimozione del Testo Inutile (Noise Removal)}
I post analizzati provengono da Reddit e sono spesso pieni di link, indirizzi email o formattazioni sfasate (come doverosi a capo multipli). 
Abbiamo eliminato tutti questi elementi tramite \textit{Espressioni Regolari (RegEx)} perché non aggiungono alcun significato utile al riconoscimento del rischio.

\begin{lstlisting}[language=Python, caption={Cancellazione di URL e Newlines}]
import re

def remove_urls_and_newlines(text):
    # Rimuove URL espliciti completi
    text = re.sub(r'http[s]?://(?:[a-zA-Z]|[0-9]|[$-_@.&+]|[!*\(\),]|(?:%[0-9a-fA-F][0-9a-fA-F]))+', '', text)
    # Comprime ritorni a capo multipli in un unico spazio
    text = re.sub(r'\n+', ' ', text)
    return text
\end{lstlisting}

\subsection{2. Lowercasing e Conservazione della Punteggiatura}
Abbiamo trasformato tutte le lettere in minuscolo (\textit{lowercasing}) per insegnare al modello che "Aiuto" e "aiuto" sono la stessa identica parola, semplificando così il vocabolario totale.

Per i modelli di elaborazione del linguaggio naturale utilizzati un tempo, l'istinto era cancellare tutta la punteggiatura. Noi abbiamo scelto invece di \textbf{mantenere i segni chiave come punti fermi, virgole, punti interrogativi ed esclamativi}. In contesti di ansia, panico o depressione, un utente che abusa di esclamazioni (es. "Aiuto!!!") esprime uno stato di forte alterazione che il modello di Intelligenza Artificiale deve essere in grado di leggere e non ignorare.

\begin{lstlisting}[language=Python, caption={Conversione in minuscolo e salvataggio punteggiatura}]
def clean_chars_and_lower(text):
    text = text.lower()
    # Manteniamo esclusivamente la punteggiatura utile (. , ! ?)
    text = re.sub(r'[^a-zA-Z0-9\s\.,!\?]', '', text)
    
    # Stacchiamo la punteggiatura per isolarla dalle parole
    text = re.sub(r'([\.,!\?])', r' \1 ', text)
    text = re.sub(r'\s+', ' ', text).strip()
    return text
\end{lstlisting}

\subsection{3. Gestione di Stop-Words e Lemmatizzazione}
Le parole di collegamento della lingua inglese (come articoli "the", pronomi "it", congiunzioni "and") occupavano oltre il 40\% di tutto il testo, distogliendo l'attenzione del modello. Tramite la libreria \texttt{NLTK} abbiamo eliminato gran parte di queste "Stop-Words".

Tuttavia, abbiamo dovuto fare estrema attenzione alle \textbf{negazioni}. Parole come "not", "don't" o "never" sono state escluse dalla lista di rimozione e conservate forzatamente nel testo. Se le avessimo rimosse, la frase "I am not suicidal" si sarebbe trasformata in "I am suicidal", scatenando un allarme falso e pericoloso per la piattaforma. 
Successivamente abbiamo applicato la \textbf{Lemmatizzazione}, un processo che riporta i verbi e i vocaboli alla loro radice del dizionario (ad esempio "depressed", "depressing" vengono convertiti nel concetto unico "depress").

\begin{lstlisting}[language=Python, caption={Epurazione parole inutili e applicazione del Lemma}]
from nltk.corpus import stopwords
from nltk.stem import WordNetLemmatizer

default_stopwords = set(stopwords.words('english'))
# Lista di eccezioni intoccabili per via del loro significato clinico
words_to_keep = {'not', 'no', 'nor', 'don', "don't", "isn't", "wasn't", 'never'}
custom_stopwords = default_stopwords - words_to_keep

lemmatizer = WordNetLemmatizer()

def remove_stops_and_lemmatize(text):
    words = text.split()
    cleaned_words = []
    
    for word in words:
        if word in ['.', ',', '!', '?']:
            cleaned_words.append(word)
        elif word not in custom_stopwords: # Se non e' una parola vietata
            cleaned_words.append(lemmatizer.lemmatize(word))
            
    return ' '.join(cleaned_words)
\end{lstlisting}

\subsection{4. Adeguamento al Limite Architetturale (Max 2000 caratteri)}
Il dataset originale su Kaggle ospitava post lunghissimi, simili a romanzi. Purtroppo, l'applicativo reale SINTONIA, all'interno del quale il paziente compila il suo diario, possiede un limite di immissione fissato a 2.000 caratteri per evitare l'abuso dello strumento. 
Per allineare il comportamento durante l'addestramento dell'Intelligenza Artificiale ai veri dati che vedrà in produzione, abbiamo imposto una rigorosa troncatura a 2000 caratteri a tutti i documenti processati. Sono stati volutamente \textbf{preservati i frammenti molto brevi} (es. frasi da 2-3 parole come "I'm done"), poiché i comportamenti perentori e in apparenza sbrigativi costituiscono un comprovato segnale psicologico di allarme imminente.

\begin{lstlisting}[language=Python, caption={Troncatura in base ai parametri di backend SINTONIA}]
def truncate_to_sintonia_limit(text):
    if len(str(text)) > 2000:
        return str(text)[:2000]
    return str(text)
\end{lstlisting}

\subsection{5. Gestione del Language Gap (Zero-Shot Cross-Lingual Transfer)}
Un enorme divario tra questa fase e la futura operatività è la lingua: il nostro dataset di partenza è in Inglese, ma gli utenti campani della piattaforma SINTONIA scriveranno i diari in Italiano. 
Solitamente in questi casi si tenta di tradurre interamente il dataset iniziale (Machine Translation), pratica che degrada la qualità dei dati, annacqua i modi di dire dello "slang" depressivo (ad es. "feeling blue") e risulta temporalmente onerosa in fase di caricamento ed addestramento. 

La soluzione progettata per il modello SINTON-IA adotta l'approccio \textbf{Zero-Shot Cross-Lingual Transfer}.
L'idea consiste nell'utilizzare come "cervello di base" un modello Transformer Multilingua aspecifico preparato in precedenza da ricercatori terzi (ad es. XLM-RoBERTa), che è capace di capire 100 lingue diverse in contemporanea. Noi addestreremo l'Intelligenza Artificiale sfruttando nativamente i post in inglese. La logica pregressa del Transformer sarà in grado di riconoscere i "concetti appresi" (il comportamento suicidario) e replicarne in autonomia l'identificazione anche quando, una volta installato in reparto, l'input live del paziente proverrà tradotto sintatticamente in lingua italiana.

% ─────────────────────────────────────────────────────────────────────────────
\clearpage
\section{Modello Depression Prediction}
\label{sec:prep_depression}

A seguito della generazione sintetica descritta nel Capitolo~\ref{ch:data_understanding}, il dataset grezzo --- composto da 1.400.000 record (100.000 pazienti $\times$ 14 giorni) --- richiede una fase di trasformazione prima di poter essere utilizzato per l'addestramento di un modello di \textit{Machine Learning}. L'obiettivo di questa fase è condensare la serie temporale di ciascun paziente in un vettore compatto di variabili numeriche (\textit{features}) che catturino i pattern clinici associati alla depressione.

Tutte le operazioni sono documentate e riproducibili nel notebook \texttt{03\_\allowbreak{}Feature\_\allowbreak{}Engineering\_\allowbreak{}and\_\allowbreak{}Prep.ipynb}.

\subsection{Feature Engineering Temporale}
\label{subsec:feature_eng_depression}

Per ciascun paziente, i 14 giorni di check-in emotivi vengono sintetizzati in un vettore di \textbf{8 feature indipendenti}. Ogni feature è stata progettata per catturare un aspetto specifico delle dinamiche emotive che la letteratura scientifica associa alla depressione:

\begin{enumerate}
    \item \textbf{\texttt{valence\_mean}}: Media della valenza emotiva sui 14 giorni. Quantifica il tono affettivo complessivo del periodo di osservazione. Un valore mediamente negativo è il primo indicatore di disagio psicologico \cite{kuppens2010inertia}.

    \item \textbf{\texttt{valence\_std}}: Deviazione standard della valenza. Misura la \textit{volatilità emotiva}: un'elevata instabilità è associata a ridotto benessere psicologico, come dimostrato dalla meta-analisi di Houben et al. \cite{houben2015emotion}.

    \item \textbf{\texttt{valence\_ema\_3d}}: Media Mobile Esponenziale (\textit{Exponential Moving Average}) calcolata sugli ultimi 3 punti temporali. Attribuisce peso preponderante alle osservazioni più recenti, catturando il \textit{recency bias} tipico dei questionari clinici a self-report.

    \item \textbf{\texttt{valence\_trend\_5d}}: Pendenza lineare della valenza negli ultimi 5 giorni, calcolata tramite regressione polinomiale di primo grado. Isola la \textit{direzione del cambiamento emotivo}: un trend discendente può segnalare il \textit{critical slowing down} descritto da Wichers et al. \cite{wichers2016critical} come precursore di transizioni depressive.

    \item \textbf{\texttt{max\_neg\_streak}}: Numero massimo di giorni consecutivi in cui la valenza è rimasta negativa ($< 0$). Quantifica la \textit{persistenza} degli stati depressivi, un fenomeno documentato da aan het Rot et al. \cite{aanhetrot2012ema} come caratteristica centrale della depressione nel monitoraggio EMA.

    \item \textbf{\texttt{missing\_ratio}}: Proporzione di giorni in cui il paziente non ha effettuato il check-in. Costituisce un predittore informativo (\textit{informative missingness}): come documentato da Mohr et al. \cite{mohr2017personal}, il disengagement dall'app è di per sé un segnale clinico associato all'anedonia e alla fatica.

    \item \textbf{\texttt{intensity\_mean}}: Media dell'intensità emotiva riportata (scala 1-10). Cattura il fenomeno dell'\textit{appiattimento affettivo}: Rottenberg et al. \cite{rottenberg2005mood} hanno dimostrato che i soggetti depressi sperimentano emozioni meno intense sia in contesti positivi che negativi.

    \item \textbf{\texttt{dominant\_mood\_valence}}: Valenza associata allo stato d'animo più frequente (la moda) nel periodo di osservazione, mappata tramite le coordinate del Modello Circomplesso \cite{russell1980circumplex}. Indica il ``quadrante emotivo dominante'' del paziente.
\end{enumerate}

\subsection{Analisi delle Correlazioni}
\label{subsec:correlazioni_depression}

Le 8 feature estratte sono state analizzate tramite matrice di correlazione di Pearson (Figura~\ref{fig:feature_corr}) per verificare due proprietà fondamentali: (i) che ciascuna feature presenti una correlazione significativa con il target PHQ-9, confermandone la rilevanza predittiva; (ii) che nessuna feature raggiunga una correlazione così elevata ($> 0.95$) da configurare un rischio di \textit{data leakage}.

\begin{figure}[htbp]
    \centering
    \includegraphics[width=0.80\textwidth]{dp_feature_correlation.png}
    \caption{Matrice di correlazione di Pearson tra le 8 feature estratte e il punteggio PHQ-9 target.}
    \label{fig:feature_corr}
\end{figure}

\subsection{Standardizzazione e Definizione del Target}
\label{subsec:scaling_depression}

Le 8 feature presentano ordini di grandezza eterogenei: la deviazione standard della valenza produce valori decimali, mentre il conteggio delle streak negative è un intero che può raggiungere 14. Per garantire che tutti i predittori contribuiscano equamente al processo di apprendimento, le feature vengono standardizzate tramite \texttt{StandardScaler} della libreria \textit{scikit-learn}, applicando la trasformazione $Z$-score:
\[
z_i = \frac{x_i - \mu}{\sigma}
\]
che impone a ciascuna colonna una media $\mu = 0$ e una deviazione standard $\sigma = 1$. Questa normalizzazione è particolarmente importante per algoritmi sensibili alla scala dei dati, come le reti neurali e le macchine a vettori di supporto.

La \textbf{variabile target} è il punteggio PHQ-9 continuo (0--27), formulato come problema di \textbf{regressione}. A differenza di un approccio di classificazione binaria (``depresso'' vs ``non depresso''), la regressione consente una granularità clinica superiore: il modello può rilevare peggioramenti graduali (ad esempio, da score 8 a score 15) prima che il paziente attraversi una soglia critica, permettendo interventi preventivi tempestivi.

Il dataset finale, contenente le 8 feature standardizzate e il target PHQ-9 in un unico file (\texttt{depression\_\allowbreak{}model\_\allowbreak{}dataset.csv}), è pronto per l'addestramento dei modelli di \textit{Machine Learning} (Capitolo 5).

% ─────────────────────────────────────────────────────────────────────────────
\clearpage
\section{Modello Churn Prediction (Log Dati Aziendali)}
\label{sec:prep_churn}
\textit{Nota: La stesura di questa sezione sarà a cura degli altri membri del team. Riguarderà manipolazione scalare e codifica numerica in vettori categoriali dei formati estrapolati in formato tabellare CSV/JSON strutturato SINTONIA.}
